% AA-10 | Algebra abstrakcyjna | Iteracja V
\chapter[Podzielność i dziedziny ideałów głównych]{Podzielność, dziedziny euklidesowe i dziedziny ideałów głównych}
\label{aa:ch:podzielnosc}
\setlength{\parskip}{2.8pt}

\section{Od dzielenia liczb do struktury ideałów}
W pierścieniu $\mathbb Z$ możemy dzielić z resztą, wyznaczać
największy wspólny dzielnik, rozkładać liczby na czynniki
pierwsze i rozwiązywać równania liniowe z całkowitymi
współczynnikami. Okazuje się, że za tymi zjawiskami stoją
\emph{różne poziomy założeń}. Dzielenie z resztą jest
silniejszą własnością niż samo istnienie największych wspólnych
dzielników; z kolei jednoznaczność rozkładu na czynniki
nie wystarcza, aby każdy ideał był generowany przez jeden element.
Rozdział wyjaśnia, które wnioski otrzymujemy na każdym poziomie.

W całym rozdziale, jeśli nie wskazano inaczej, $R$ jest
dziedziną całkowitości z jedynką, czyli niezerowym pierścieniem
przemiennym bez dzielników zera. Założenie to jest ważne:
pozwala skracać niezerowe czynniki i porównywać dzielniki.

\begin{definicja}[Podzielność i elementy stowarzyszone]
\label{aa:def:divides-associates}
Niech $R$ będzie dziedziną całkowitości z jedynką.
Dla $a,b\in R$ zapis $a\mid b$ oznacza, że istnieje
$c\in R$ z $b=ac$. Elementy $a,b$ są \emph{stowarzyszone},
jeśli $a=ub$ dla pewnej jednostki $u\in R^\times$.
Dla $a,b\in R$ \emph{największym wspólnym dzielnikiem}
nazywamy element $d\in R$ taki, że $d\mid a$, $d\mid b$
oraz każdy wspólny dzielnik $a,b$ dzieli $d$.
Jeżeli taki element istnieje, oznaczamy go przez
$\gcd(a,b)$, pamiętając, że jest wyznaczony tylko
z dokładnością do stowarzyszenia.
\end{definicja}
\begin{lemat}[Stowarzyszenie i jednoznaczność największego wspólnego dzielnika]
\label{aa:lem:gcd-unique}
Niech $R$ będzie dziedziną całkowitości. Elementy
$a,b\in R\setminus\{0\}$ są stowarzyszone wtedy i tylko
wtedy, gdy $a\mid b$ oraz $b\mid a$.
Dla dowolnych $a,b\in R$, jeżeli $d,d'\in R$ są największymi
wspólnymi dzielnikami pary $a,b$, to są stowarzyszone.
\end{lemat}
\begin{proof}
Z $a=ub$ i odwracalności $u$ natychmiast wynika
$a\mid b$ i $b\mid a$. W drugą stronę, jeśli $a=bc$
i $b=ad$, to $a=acd$. Ponieważ $a\ne0$, skracanie
daje $cd=1_R$, a więc $c$ i $d$ są jednostkami.
Dla dwóch największych wspólnych dzielników $d,d'$
warunek ich maksymalności w sensie podzielności
daje $d\mid d'$ i $d'\mid d$. Jeśli $a=b=0$, to każdy element dzieli zarówno $a$, jak i $b$.
W szczególności $0$ jest wspólnym dzielnikiem, więc z definicji
największego wspólnego dzielnika każdy taki element $d$ musi
spełniać $0\mid d$, czyli $d=0$. Element $0$ sam spełnia
warunek największego wspólnego dzielnika, ponieważ każdy element
dzieli $0$. Zatem $\gcd(0,0)=0$ także w sensie definicji.
W pozostałych przypadkach $d,d'$ są niezerowe i stosujemy
pierwszą część.
\end{proof}

\begin{przyklad}[Jednostki decydują o jednoznaczności]
\label{aa:prz:units-gcd}
W $\mathbb Z$ elementy stowarzyszone z $6$ to $6$ oraz $-6$;
wybieramy dodatni największy wspólny dzielnik.
W pierścieniu $\mathbb Q[x]$ elementy $x^2-1$ oraz
$7(x^2-1)$ są stowarzyszone, ponieważ $7$ jest jednostką.
Największy wspólny dzielnik wielomianów zwykle wybiera się
\emph{unormowany}, to znaczy o najwyższym współczynniku
równym $1$, jeśli wielomian nie jest zerowy.
\end{przyklad}

\begin{definicja}[Element nierozkładalny i element pierwszy]
\label{aa:def:irreducible-prime}
Niech $R$ będzie dziedziną całkowitości z jedynką i niech
$p\in R$ będzie niezerowy i nie będzie jednostką.
Element $p$ jest \emph{nierozkładalny}, jeśli z $p=ab$
wynika, że $a$ lub $b$ jest jednostką.
Element $p$ jest \emph{pierwszy}, jeśli dla wszystkich
$a,b\in R$ z $p\mid ab$ wynika $p\mid a$ lub $p\mid b$.
\end{definicja}
\begin{lemat}[Pierwszość implikuje nierozkładalność]
\label{aa:lem:prime-irred}
Niech $R$ będzie dziedziną całkowitości z jedynką.
Każdy element pierwszy $p\in R$ jest nierozkładalny.
\end{lemat}
\begin{proof}
Niech $p=ab$. Wtedy $p\mid ab$, a pierwszość
pozwala przyjąć, po ewentualnej zamianie $a,b$, że $p\mid a$.
Istnieje $c$ z $a=pc$, zatem $p=pcb$.
Skracając $p\ne0$, otrzymujemy $1_R=cb$.
Element $b$ jest jednostką. To dokładnie definicja nierozkładalności.
\end{proof}

\section{Dziedziny euklidesowe}
\begin{definicja}[Dziedzina euklidesowa]
\label{aa:def:euclidean-domain}
\emph{Dziedziną euklidesową} nazywamy dziedzinę całkowitości
$R$, dla której istnieje funkcja
$\delta:R\setminus\{0\}\to\mathbb N_0$ taka, że dla
każdych $a\in R$ i $b\in R\setminus\{0\}$ istnieją
$q,r\in R$ spełniające $a=bq+r$ oraz
$r=0$ albo $\delta(r)<\delta(b)$.
Funkcję $\delta$ nazywamy \emph{funkcją euklidesową}.
Nie zakładamy jednoznaczności $q,r$.
\end{definicja}

\begin{przyklad}[Liczby całkowite i dzielenie z resztą]
\label{aa:prz:integer-division}
Dla $R=\mathbb Z$ funkcja $\delta(b)=|b|$, określona dla
$b\in\mathbb Z\setminus\{0\}$, jest euklidesowa.
Jeżeli $b>0$, dzielenie z resztą daje jednoznaczne
$a=bq+r$ przy $0\le r<b$.
Gdy $b<0$, przechodzimy do $|b|$ i zmieniamy znak ilorazu.
Dla $a=-17$, $b=5$ otrzymujemy
$-17=5(-4)+3$, a $0\le3<5$. Wybór $q=-3$ dałby resztę
$-2$, która też spełnia $|{-2}|<5$, ale nie standardowy
warunek $0\le r<5$. Dlatego funkcja euklidesowa nie narzuca
jednoznaczności dzielenia, jeżeli nie ustalimy dodatkowej
konwencji dla reszt.
\end{przyklad}

\begin{przyklad}[Wielomiany jednej zmiennej]
\label{aa:prz:polynomial-division}
Niech $K$ będzie ciałem. W $K[x]$ funkcja $\delta(f)=\deg f$
dla $f\ne0$ jest euklidesowa.
Wyjaśnijmy mechanizm. Jeśli $f=0$ lub $\deg f<\deg g$, wystarczy wziąć
$q=0$ i $r=f$. W pozostałym przypadku $f\ne0$ oraz $\deg f\ge\deg g$.
Niech $a,b\in K^\times$ będą najwyższymi współczynnikami
odpowiednio $f$ i $g$. Odejmujemy od $f$ wielomian
$(a/b)x^{\deg f-\deg g}g$.
Najwyższy wyraz znika; nowy wielomian ma mniejszy stopień.
Powtarzanie kroku kończy się po skończonej liczbie operacji,
bo stopnie nie mogą maleć w nieskończoność.
Dostajemy $f=gq+r$, gdzie $r=0$ lub $\deg r<\deg g$.
Na przykład nad $\mathbb Q$:
\[
x^3+2x+1=(x+1)(x^2-x+3)-2.
\]
Rachunek można sprawdzić przez wymnożenie:
$(x+1)(x^2-x+3)=x^3+2x+3$.
Stała reszta $-2$ ma stopień $0<1$, więc spełnia wymóg.
\end{przyklad}

\begin{przyklad}[Liczby całkowite Gaussa]
\label{aa:prz:gaussian-euclidean}
Pierścień $\mathbb Z[i]=\{a+bi:a,b\in\mathbb Z\}$
jest dziedziną euklidesową względem funkcji euklidesowej
$\delta(z)=N(z)$ dla $z\ne0$, gdzie
$N(a+bi)=a^2+b^2$ jest określona dla wszystkich
$a+bi\in\mathbb Z[i]$, w tym dla zera.
Jeśli $\alpha,\beta\in\mathbb Z[i]$ i $\beta\ne0$,
zapiszmy $\alpha/\beta=x+iy$ z $x,y\in\mathbb R$.
Wybierzmy liczby całkowite $m,n$ najbliższe odpowiednio
$x,y$; wtedy $|x-m|\le\tfrac12$ i $|y-n|\le\tfrac12$.
Dla $q=m+in$ oraz $r=\alpha-\beta q$ zachodzi
$\alpha=\beta q+r$ i
\[
N(r)=N(\beta)\,N(\alpha/\beta-q)
\le N(\beta)(1/4+1/4)=N(\beta)/2<N(\beta).
\]
Ponieważ $r\in\mathbb Z[i]$, uzyskaliśmy wymagane dzielenie
z resztą. Własność multiplikatywną normy można sprawdzić
bez odwołania do analizy: dla $z=a+bi$, $w=c+di$
mamy $zw=(ac-bd)+(ad+bc)i$, a zatem
\[
N(zw)=(ac-bd)^2+(ad+bc)^2
=(a^2+b^2)(c^2+d^2)=N(z)N(w).
\]
Tutaj $i^2=-1$; konstrukcję ciała liczb zespolonych
z danej struktury ciała liczb rzeczywistych podamy
w dalszym rozdziale o konstrukcjach algebraicznych.
\end{przyklad}

\section{Algorytm Euklidesa i tożsamość Bézouta}
\begin{twierdzenie}[Algorytm Euklidesa]
\label{aa:tw:euclidean-algorithm}
Niech $R$ będzie dziedziną euklidesową z funkcją euklidesową
$\delta$ i niech $a,b\in R$, $b\ne0$. Połóżmy
$r_{-1}=a$ oraz $r_0=b$. Istnieje liczba całkowita $k\ge0$ i istnieją elementy
$q_0,\ldots,q_k,r_1,\ldots,r_{k+1}\in R$
takie, że
\[
r_{j-1}=r_jq_j+r_{j+1}\quad(0\le j\le k),
\qquad r_{k+1}=0,
\]
$r_j\ne0$ dla $0\le j\le k$ oraz
$\delta(r_{j+1})<\delta(r_j)$ dla $0\le j<k$.
Ostatni niezerowy element $r_k$ jest największym wspólnym
dzielnikiem $a$ i $b$. W szczególności, gdy już pierwsza
reszta jest zerem, mamy $k=0$ i $r_k=b$.
Co więcej, istnieją $u,v\in R$ takie, że
$ua+vb=\gcd(a,b)$, z dokładnością do wyboru elementu
stowarzyszonego z największym wspólnym dzielnikiem.
\end{twierdzenie}
\begin{proof}
Dzieląc kolejne pary elementów, w każdym kroku otrzymujemy
resztę równą zeru albo o mniejszej wartości $\delta$.
Wartości $\delta$ należą do $\mathbb N_0$, więc nie istnieje
nieskończony ściśle malejący ciąg wartości. Algorytm się kończy.
Z równości $a=bq_0+r_1$ wynika, że element $d$ dzieli
jednocześnie $a,b$ wtedy i tylko wtedy, gdy dzieli $b,r_1$:
jeśli $d\mid a,b$, to $d\mid a-bq_0=r_1$;
jeśli $d\mid b,r_1$, to $d\mid bq_0+r_1=a$.
Taką samą równoważność stosujemy na każdym kolejnym kroku.
Jeżeli pierwsza reszta jest zerem, ostatnim niezerowym
elementem jest $b$, którego dzielniki są wspólnymi dzielnikami
$a$ i $b$. W przeciwnym razie wspólni dzielnicy $a,b$ są
dokładnie dzielnikami ostatniej niezerowej reszty $r_k$: dzielnik $r_k$
dzieli także poprzednią resztę $r_{k-1}=r_k q_k$,
a wszystkie wcześniejsze na podstawie odwróconych równości.
To dowodzi charakterystyki największego wspólnego dzielnika.

Aby uzyskać tożsamość Bézouta, wykonujemy podstawienia wsteczne.
Pierwsza reszta ma postać $r_1=a-bq_0$,
następna $r_2=b-r_1q_1$, a więc też jest kombinacją
$R$-liniową $a,b$. Indukcyjnie każda kolejna reszta
jest kombinacją liniową tych samych dwóch elementów.
W przypadku zerowej pierwszej reszty mamy $b=0a+1b$.
W przeciwnym razie $r_k=ua+vb$ dla pewnych $u,v\in R$.
\end{proof}

\begin{przyklad}[Rozszerzony algorytm Euklidesa dla $252$ i $198$]
\label{aa:prz:extended-euclid}
W $\mathbb Z$ wykonujemy dzielenie:
\[
252=1\cdot198+54,\qquad
198=3\cdot54+36,\qquad
54=1\cdot36+18,\qquad36=2\cdot18.
\]
Ostatnia niezerowa reszta wynosi $18$.
Podstawiamy od końca:
\[
\begin{aligned}
18&=54-36=54-(198-3\cdot54)\\
  &=4\cdot54-198=4(252-198)-198\\
  &=4\cdot252-5\cdot198.
\end{aligned}
\]
Otrzymujemy nie tylko $\gcd(252,198)=18$, lecz także
konkretne współczynniki Bézouta. Ponieważ $18$ jest
największym wspólnym dzielnikiem, równanie
$252x+198y=18$ ma rozwiązanie całkowite $(x,y)=(4,-5)$.
\end{przyklad}

\begin{wniosek}[Lemat Euklidesa w dziedzinie euklidesowej]
\label{aa:wn:euclid-lemma}
Niech $a,b,c$ należą do dziedziny euklidesowej $R$.
Jeżeli $\gcd(a,b)$ jest jednostką oraz $a\mid bc$, to $a\mid c$.
\end{wniosek}
\begin{proof}
Tożsamość Bézouta daje $ua+vb=1_R$ dla pewnych $u,v\in R$
(po pomnożeniu przez odwrotność największego wspólnego dzielnika).
Mnożąc przez $c$, otrzymujemy $c=uac+vbc$.
Pierwszy składnik jest podzielny przez $a$ z definicji,
a drugi na mocy $a\mid bc$. Zatem $a\mid c$.
\end{proof}

\begin{przyklad}[Największy wspólny dzielnik wielomianów]
\label{aa:prz:poly-gcd}
W $\mathbb Q[x]$ rozważmy $f=x^3-1$ i $g=x^2-1$.
Mamy $f=(x-1)(x^2+x+1)$ oraz $g=(x-1)(x+1)$.
Dzielenie $x^2+x+1$ przez $x+1$ daje resztę $1$,
bo podstawienie $x=-1$ daje $1$. Zatem oba czynniki
$x^2+x+1$ i $x+1$ są względnie pierwsze.
Tożsamość Bézouta dla tych dwóch wielomianów daje
$u(x)(x^2+x+1)+v(x)(x+1)=1$ dla pewnych $u,v\in\mathbb Q[x]$.
Po pomnożeniu przez $x-1$ otrzymujemy
$u(x)f(x)+v(x)g(x)=x-1$. Każdy wspólny dzielnik
$f$ i $g$ dzieli więc $x-1$, który sam dzieli oba wielomiany.
Otrzymujemy unormowany największy wspólny dzielnik
$\gcd(f,g)=x-1$.
\end{przyklad}

\section{Dziedziny ideałów głównych}
\begin{definicja}[Dziedzina ideałów głównych]
\label{aa:def:pid}
Dziedzinę całkowitości $R$, w której każdy ideał ma postać
$(d)=\{rd:r\in R\}$ dla pewnego $d\in R$, nazywamy
\emph{dziedziną ideałów głównych}. Oznacza to, że nawet
ideał generowany przez wiele elementów można wygenerować
jednym elementem. Nie znaczy to, że wszystkie elementy
pierścienia są wzajemnie podzielne.
\end{definicja}
\begin{twierdzenie}[Każda dziedzina euklidesowa jest dziedziną ideałów głównych]
\label{aa:tw:euclidean-pid}
Jeżeli $R$ jest dziedziną euklidesową, to każdy ideał
$R$ jest główny.
\end{twierdzenie}
\begin{proof}
Niech $I$ będzie ideałem $R$. Jeśli $I=\{0\}$,
to $I=(0)$. W przeciwnym razie rozważmy niepusty zbiór
$\{\delta(a):a\in I,\ a\ne0\}\subseteq\mathbb N_0$.
Ma on element najmniejszy. Wybierzmy niezerowy $d\in I$
o minimalnej wartości $\delta(d)$. Dla dowolnego $a\in I$
wykonajmy dzielenie $a=dq+r$, gdzie $r=0$
lub $\delta(r)<\delta(d)$. Ponieważ $a,d\in I$
oraz $I$ jest ideałem, $r=a-dq\in I$.
Gdyby $r\ne0$, przeczyłoby to minimalności $\delta(d)$.
Zatem $r=0$, czyli $a=dq\in(d)$. Udowodniliśmy $I\subseteq(d)$.
Odwrotna inkluzja zachodzi, bo $d\in I$ i $I$ jest ideałem.
Stąd $I=(d)$.
\end{proof}

\begin{twierdzenie}[Największy wspólny dzielnik w dziedzinie ideałów głównych]
\label{aa:tw:pid-bezout}
Jeżeli $R$ jest dziedziną ideałów głównych, to dla
każdych $a,b\in R$, z których co najmniej jeden jest niezerowy,
ideał $(a,b)$ ma postać $(d)$, gdzie $d$ jest największym
wspólnym dzielnikiem $a,b$. Istnieją $u,v\in R$ z
$d=ua+vb$.
\end{twierdzenie}
\begin{proof}
Z definicji dziedziny ideałów głównych istnieje $d\in R$
z $(a,b)=(d)$. Ponieważ $a,b\in(d)$, mamy $d\mid a$
i $d\mid b$. Z drugiej strony $d\in(a,b)$, więc
dla pewnych $u,v\in R$ zachodzi $d=ua+vb$.
Jeśli $c$ jest wspólnym dzielnikiem $a,b$, to dzieli każdą
ich kombinację liniową, a zatem dzieli $d$.
W ten sposób wykazaliśmy zarówno własność największego
wspólnego dzielnika, jak i tożsamość Bézouta.
\end{proof}

\begin{wniosek}[Pierwszość elementu nierozkładalnego w dziedzinie ideałów głównych]
\label{aa:wn:irred-is-prime-pid}
Niech $R$ będzie dziedziną ideałów głównych i niech
$p\in R$ będzie elementem nierozkładalnym. Wtedy $p$ jest pierwszy.
\end{wniosek}
\begin{proof}
Niech $p$ będzie nierozkładalny i niech $p\mid ab$.
Jeśli $p\nmid a$, to największy wspólny dzielnik $p,a$
musi być jednostką: dzieli $p$, a nierozkładalność $p$
oznacza, że każdy jego dzielnik jest jednostką albo jest
stowarzyszony z $p$; druga możliwość wymuszałaby $p\mid a$.
Tożsamość Bézouta daje $up+va=1_R$.
Mnożąc przez $b$, otrzymujemy $b=upb+vab$;
oba składniki są podzielne przez $p$. A zatem $p\mid b$.
\end{proof}

\section{Dlaczego rozkład na czynniki musi się kończyć?}
\begin{definicja}[Dziedzina z jednoznacznym rozkładem]
\label{aa:def:ufd}
Dziedzinę całkowitości $R$ nazywamy
\emph{dziedziną z jednoznacznym rozkładem na czynniki},
jeśli każdy niezerowy element niebędący jednostką można
zapisać jako skończony iloczyn elementów nierozkładalnych,
a dwa takie zapisy różnią się co najwyżej kolejnością
czynników i zastąpieniem każdego czynnika elementem stowarzyszonym.
\end{definicja}

\begin{lemat}[Warunek stabilizacji wstępujących ciągów ideałów]
\label{aa:lem:pid-ascending-chain}
Niech $R$ będzie dziedziną ideałów głównych.
Każdy wstępujący ciąg ideałów $R$,
$I_1\subseteq I_2\subseteq\cdots$, staje się stały.
\end{lemat}
\begin{proof}
Suma zbiorów $I=\bigcup_{n\ge1}I_n$ jest ideałem:
jeżeli $a\in I_j$ i $b\in I_k$, to obydwa należą
do $I_{\max(j,k)}$, więc $a-b$ należy do tego ideału;
to samo dotyczy mnożenia przez dowolny element pierścienia.
Ponieważ $R$ jest dziedziną ideałów głównych, $I=(d)$.
Generator $d$ należy do $I_N$ dla pewnego $N$,
zatem $(d)\subseteq I_N$. Z drugiej strony $I_N\subseteq I$.
Wynika stąd $I_N=I$, a dla $n\ge N$ mamy $I_n=I$.
\end{proof}

\begin{twierdzenie}[Dziedziny ideałów głównych mają jednoznaczny rozkład na czynniki]
\label{aa:tw:pid-ufd}
Niech $R$ będzie dziedziną ideałów głównych.
Wówczas $R$ jest dziedziną z jednoznacznym rozkładem
na czynniki.
\end{twierdzenie}
\begin{proof}
Najpierw udowodnimy istnienie rozkładu.
Przypuśćmy przeciwnie, że istnieje niezerowy nieodwracalny
element $a_0$, którego nie można przedstawić jako skończonego
iloczynu elementów nierozkładalnych. Nie jest on nierozkładalny,
bo wtedy iloczyn złożony z jednego czynnika byłby odpowiednim
rozkładem. Istnieją więc nieodwracalne $a_1,b_1$ z
$a_0=a_1b_1$. Co najmniej jeden z czynników nie ma rozkładu
na nierozkładalne: gdyby oba miały, ich rozkłady po pomnożeniu
dałyby rozkład $a_0$. Oznaczmy ten czynnik przez $a_1$
(po ewentualnej zamianie czynników).
Powtarzając argument, otrzymujemy ciąg
$a_0=a_1b_1$, $a_1=a_2b_2$, itd.,
gdzie każde $b_i$ nie jest jednostką. Stąd
$(a_0)\subseteq(a_1)\subseteq(a_2)\subseteq\cdots$.
Każde zawieranie jest ścisłe: gdyby $(a_{j-1})=(a_j)$,
to dla pewnego $d\in R$ mielibyśmy $a_j=a_{j-1}d$.
Z równości $a_{j-1}=a_jb_j$ wynikałoby wtedy
$a_j=a_jb_jd$. Ponieważ $a_j\ne0$, skracanie daje
$b_jd=1_R$, czyli $b_j$ jest jednostką, wbrew konstrukcji.
Otrzymujemy nieskończony ściśle rosnący ciąg ideałów,
sprzeczny z poprzednim lematem. Zatem rozkład istnieje.

Teraz udowodnimy jednoznaczność. Każdy element nierozkładalny
jest pierwszy na mocy Wniosku~\ref{aa:wn:irred-is-prime-pid}.
Jeżeli $p_1\cdots p_r=q_1\cdots q_s$ są dwoma rozkładami,
to $p_1$ dzieli iloczyn po prawej. Wielokrotne zastosowanie
pierwszości daje, że $p_1\mid q_j$ dla pewnego $j$.
Ponieważ $q_j$ jest nierozkładalny, $p_1$ i $q_j$
są stowarzyszone. Po przestawieniu czynników możemy
założyć $j=1$ i zapisać $q_1=up_1$ dla jednostki $u$.
Skracając niezerowe $p_1$, otrzymujemy równość iloczynu
pozostałych $p_i$ z iloczynem pozostałych $q_j$ pomnożonym
przez jednostkę $u$. Jeżeli po obu stronach są jeszcze czynniki,
wchłaniamy $u$ w jeden z nich; nie zmienia to nierozkładalności.
Przypadek, gdy po obu stronach nie pozostał żaden czynnik,
kończy indukcję. Gdy po skróceniu pozostałyby
po jednej stronie czynniki nierozkładalne, a po drugiej
jednostka, powstałaby sprzeczność: iloczyn niejednostek
w dziedzinie całkowitości nie może być jednostką.
Indukcja daje $r=s$ oraz stowarzyszenie odpowiadających sobie
czynników. Gdy jeden z rozkładów jest pusty, obie strony
są jednostkami i nie mogą jednocześnie zawierać czynników
nierozkładalnych. To kończy dowód.
\end{proof}

\begin{uwaga}[Łańcuch implikacji i pytania odwrotne]
\label{aa:uw:euclidean-pid-ufd}
Udowodniliśmy implikacje
\[
\begin{gathered}
\text{dziedzina euklidesowa}\ \Longrightarrow\
\text{dziedzina ideałów głównych}\\
\Longrightarrow\ \text{dziedzina z jednoznacznym rozkładem na czynniki}.
\end{gathered}
\]
Odwrotne implikacje nie wynikają z powyższych dowodów.
Nie każdy pierścień z jednoznacznym rozkładem na czynniki
ma wszystkie ideały główne. Przykład taki uzyskamy przy
pierścieniu $K[x,y]$, gdy udowodnimy jego własność
jednoznacznego rozkładu w następnym rozdziale.
Istnieją również dziedziny ideałów głównych, które nie są
euklidesowe, lecz ich konstrukcja i dowód tego faktu
wykraczają poza obecny rozdział.
\end{uwaga}

\begin{przyklad}[Dlaczego pierścień wielomianów dwóch zmiennych nie jest dziedziną ideałów głównych]
\label{aa:prz:two-polynomials-not-pid}
Niech $K$ będzie dowolnym ciałem. W pierścieniu $K[x,y]$
rozważmy ideał $I=(x,y)$. Załóżmy, że $I=(f)$.
Wtedy $f$ dzieli jednocześnie $x$ i $y$.
Ponieważ $K$ jest ciałem, pierścień $K[x,y]$ jest
dziedziną całkowitości. Rzeczywiście, traktując niezerowe wielomiany $f,g$ jako
wielomiany zmiennej $y$ o współczynnikach w $K[x]$,
ich najwyższe współczynniki względem $y$ są niezerowymi
wielomianami w $x$. Ich iloczyn jest niezerowy, gdyż
iloczyn dwóch niezerowych wielomianów jednej zmiennej
nad ciałem ma niezerowy najwyższy współczynnik.
To dowodzi, że $K[x,y]$ nie ma dzielników zera.
Ponadto stopień całkowity niezerowego iloczynu
jest sumą stopni całkowitych: jeśli $f_d$ i $g_e$
są niezerowymi częściami jednorodnymi najwyższych
stopni całkowitych $d$ i $e$, to częścią jednorodną
stopnia $d+e$ iloczynu jest $f_dg_e$, która jest
niezerowa na mocy właśnie wykazanej własności. Z podzielności $f\mid x$ wynika,
że $f$ ma stopień najwyżej $1$. Jeśli ma stopień $1$,
to musi być stowarzyszony z $x$; analogicznie z $f\mid y$
wynikałoby stowarzyszenie z $y$. Nie jest to możliwe, bo jednostkami
$K[x,y]$ są niezerowe stałe, a równość $x=cy$ z
$c\in K^\times$ obala podstawienie $y=0$.
Zatem $f$ ma stopień $0$, jest niezerową stałą i jest jednostką.
Wtedy $I=K[x,y]$, czyli $1\in(x,y)$, ale to sprzeczne
z podstawieniem $x=y=0$ do równości $1=xf_1+yf_2$.
Wniosek: ideał $(x,y)$ nie jest główny.
\end{przyklad}

\section{Gdy zawodzi jednoznaczny rozkład}
\begin{przyklad}[Dwa różne rozkłady liczby $6$]
\label{aa:prz:zsqrtminus5}
Rozważmy dziedzinę całkowitości
$R=\mathbb Z[\sqrt{-5}]=\{a+b\sqrt{-5}:a,b\in\mathbb Z\}$.
Jest ona podpierścieniem $\mathbb C$, więc nie ma dzielników
zera. Norma $N(a+b\sqrt{-5})=a^2+5b^2$ ma wartości w
$\mathbb N_0$ i jest multiplikatywna:
$N(zw)=N(z)N(w)$. Rzeczywiście, dla sprzężenia
$\overline{a+b\sqrt{-5}}=a-b\sqrt{-5}$ mamy
$N(z)=z\overline z$, a więc
$N(zw)=zw\overline z\,\overline w=N(z)N(w)$. W tej dziedzinie
\[
6=2\cdot3=(1+\sqrt{-5})(1-\sqrt{-5}).
\]
Wykażemy, że wszystkie cztery czynniki są nierozkładalne.
Najpierw wyznaczmy wszystkie jednostki. Jeśli $u$ ma
odwrotność w $R$, to $N(u)N(u^{-1})=1$. Obie normy są
dodatnimi liczbami całkowitymi, więc $N(u)=1$.
Równanie $a^2+5b^2=1$ wymusza $b=0$ i $a=\pm1$.
Odwrotnie, $1$ oraz $-1$ są jednostkami, bo są swoimi
własnymi odwrotnościami. Są to zatem wszystkie jednostki.
Nie istnieją elementy o normie $2$ ani $3$:
przy $b\ne0$ norma jest co najmniej $5$, a przy $b=0$
musiałaby być kwadratem liczby całkowitej.
Jeśli $2=ab$ z nieodwracalnymi $a,b$, to
$4=N(2)=N(a)N(b)$, a oba czynniki normowe byłyby
większe od $1$; jedyna możliwość to $2$ i $2$, niemożliwa.
Tak samo $N(3)=9$ wymagałoby dla rozkładu na niejednostki
norm $3$ i $3$, których nie ma. Dla
$1\pm\sqrt{-5}$ norma wynosi $6$, więc rozkład
na niejednostki wymagałby norm $2$ i $3$, również niemożliwych.
Wreszcie elementy $2,3,1+\sqrt{-5},1-\sqrt{-5}$
nie są parami stowarzyszone: mają różne normy, z wyjątkiem
ostatniej pary, ale $1+\sqrt{-5}\ne\pm(1-\sqrt{-5})$.
Otrzymaliśmy dwa rozkłady na nierozkładalne, które nie
różnią się jedynie kolejnością i jednostkami.
Pierścień $\mathbb Z[\sqrt{-5}]$ nie jest więc dziedziną
z jednoznacznym rozkładem na czynniki.
\end{przyklad}

\begin{przyklad}[Element nierozkładalny, który nie jest pierwszy]
\label{aa:prz:irred-not-prime}
W powyższym pierścieniu liczba $2$ jest nierozkładalna,
ale nie jest elementem pierwszym. Mamy
$2\mid (1+\sqrt{-5})(1-\sqrt{-5})$, ponieważ iloczyn
wynosi $6=2\cdot3$. Tymczasem $2$ nie dzieli żadnego
z czynników: równości
$1\pm\sqrt{-5}=2(a+b\sqrt{-5})$ dla $a,b\in\mathbb Z$
wymagałyby $2a=1$, co nie ma rozwiązania całkowitego.
Widać dokładnie, dlaczego implikacja
„nierozkładalny $\Rightarrow$ pierwszy” wymaga dodatkowych
założeń o pierścieniu.
\end{przyklad}

\begin{uwaga}[Związek z modułami i algebrą liniową]
\label{aa:uw:pid-module-link}
Dziedzina ideałów głównych umożliwia nie tylko jednoznaczny
rozkład elementów. Jej głębszą cechą jest to, że podmoduły
modułów wolnych skończonej rangi również są wolne i mają
kontrolowaną postać. Z tego wyprowadzimy klasyfikację
skończenie generowanych modułów nad dziedziną ideałów
głównych oraz klasyfikację skończenie generowanych grup
abelowych. Następny rozdział rozwinie pierścienie wielomianów
i kryteria nierozkładalności, potrzebne również w teorii Galois.
\end{uwaga}
\setlength{\parskip}{4pt}
